\section*{الفصل \ref{ch:b2:reaction-free-energy} --- إنتروبيا التفاعل وطاقته الحرة}

\begin{solution}{exo:b2:reaction-free-energy:1}
(أ) ثلاثة مولات من الغاز تصير سائلًا: $\Delta_r S^\circ < 0$. (ب) يتكوّن غاز:
$> 0$. (ج) جزيء غازي واحد يعطي جزيئين: $> 0$. (د) أيونات في محلول
تكوّن بلورة: $< 0$. (ه) انصهار: $> 0$. (و) $\Delta\nu_{\text{غاز}} = 0$:
قيمة صغيرة، وهي هنا موجبة (\qty{+20}{J/(K.mol)} من الجداول، لأن \ce{HCl}
أقل تناظرًا من \ce{H2} و\ce{Cl2}).
% ledger: s0:HCl_g, s0:H2_g, s0:Cl2_g
\end{solution}

\begin{solution}{exo:b2:reaction-free-energy:2}
$\Delta_r S^\circ = 2(192.77) - 191.61 - 3(130.68) = \qty{-198.1}{J/(K.mol)}$؛
$\Delta_r G^\circ = -91.9 - 298.15 \times (-0.1981) = \qty{-32.8}{kJ/mol}$:
فالتفاعل \omterm{def:b2:reaction-free-energy:exergonic}{طارد للطاقة} عند \qty{298}{K}، وإن كان حد الإنتروبيا
(\qty{+59.1}{kJ/mol}) يعمل ضده.
\end{solution}

\begin{solution}{exo:b2:reaction-free-energy:3}
$\Delta_r S^\circ = 69.95 - 130.68 - \tfrac12(205.15) = \qty{-163.3}{J/(K.mol)}$؛
$\Delta_r G^\circ = -285.83 + 298.15 \times 0.1633 = \qty{-237.14}{kJ/mol}$.
إنه تفاعل تكوين الماء السائل (مول واحد من العناصر في حالاتها المرجعية)،
\omterm{def:b2:reaction-free-energy:reaction-gibbs}{فطاقة جيبس القياسية لتفاعله} هي بالتعريف
$\Delta_f G^\circ(\ce{H2O}, \text{l})$؛ وتعطي الجداول القيمة
نفسها.
\end{solution}

\begin{solution}{exo:b2:reaction-free-energy:4}
نحو \qty{42}{J/(K.mol)} عند \qty{298}{K} و\qty{87}{J/(K.mol)} عند
\qty{1000}{K} (سائل). القفزتان: \qty{10.6}{J/(K.mol)} عند \qty{692.7}{K}
و\qty{97.7}{J/(K.mol)} عند \qty{1180.2}{K}. $\Delta_{\text{انصهار}}H = 10.6 \times
692.7 = \qty{7.3}{kJ/mol}$ و$\Delta_{\text{تبخر}}H = 97.7 \times 1180.2 =
\qty{115}{kJ/mol}$.
\end{solution}

\begin{solution}{exo:b2:reaction-free-energy:5}
$\Delta_r H^\circ = \qty{44.00}{kJ/mol}$، $\Delta_r S^\circ = 188.84 - 69.95 =
\qty{118.9}{J/(K.mol)}$، $\Delta_r G^\circ = 44.00 - 298.15 \times 0.1189 =
\qty{+8.55}{kJ/mol}$. عند \qty{298}{K} لا يكون التحوّل عكوسًا (فالماء السائل
وبخاره عند \qty{1}{bar} ليسا في توازن)، ومنه $\Delta S \ne
Q/T$: $\Delta_r H^\circ/T = \qty{147.6}{J/(K.mol)}$. ويعني كون
$\Delta_r G^\circ$ موجبًا أن الماء عند \qty{298}{K} لا يتبخر في جوّ من البخار
النقي عند \qty{1}{bar}: بل يتكثف البخار. ويكون الطوران في توازن حيث
$\Delta_r G^\circ = 0$: $T = 44\,004/118.9 =
\qty{370}{K}$، أي في حدود \qty{3}{K} من نقطة الغليان النظامية
(\qty{373}{K}): \omterm{def:b2:reaction-free-energy:ellingham-approximation}{فتقريب إلينغهام} جيد على \qty{75}{K}.
\end{solution}

\begin{solution}{exo:b2:reaction-free-energy:6}
$T_i = 57\,100/175.7 = \qty{325}{K}$. وتحتها يكون $\Delta_r G^\circ > 0$
ويُرجَّح \ce{N2O4} عديم اللون؛ فتبريد الأنبوب يزيح الغاز نحو
\ce{N2O4} ويبهت اللون البني لثنائي أكسيد النيتروجين \ce{NO2} (والصلة الكمية بالتركيب
في \cref{ch:b2:equilibrium-shifts}).
\end{solution}

\begin{solution}{exo:b2:reaction-free-energy:7}
$\Delta_r S^\circ = -\dd\Delta_r G^\circ/\dd T = -(-4.0 + 12.0)/100 =
\qty{-0.080}{kJ/(K.mol)} = \qty{-80}{J/(K.mol)}$ و$\Delta_r H^\circ =
\Delta_r G^\circ + T\Delta_r S^\circ = -12.0 + 300 \times (-0.080) =
\qty{-36.0}{kJ/mol}$. التحقق: تنتقل $\Delta_r G^\circ/T$ من $-0.0400$ إلى
$\qty{-0.0100}{kJ/(K.mol)}$، بميل $3.00 \times 10^{-4}$ لكل كلفن، يساوي
$-\Delta_r H^\circ/T^2 = 36.0/346^2 = 3.00 \times 10^{-4}$ عند درجة الحرارة
المتوسطة الهندسية $\sqrt{300 \times 400} = \qty{346}{K}$.
\end{solution}

\begin{solution}{exo:b2:reaction-free-energy:8}
$\Delta_r S^\circ = 186.25 - 5.74 - 2(130.68) = \qty{-80.8}{J/(K.mol)}$،
$\Delta_r G^\circ = -74.87 + 298.15 \times 0.0808 = \qty{-50.8}{kJ/mol}$،
وهي قيمة $\Delta_f G^\circ$ المجدولة للميثان. ولأن للحدين الإشارة
نفسها، يصير الميثان غير مستقر فوق $T_i = 74\,870/80.8 \approx
\qty{926}{K}$ (\omterm{def:b2:reaction-free-energy:ellingham-approximation}{تقريب إلينغهام}): فالميثان الساخن يتكسر إلى كربون
وهيدروجين، وهي طريق لإنتاجهما كليهما.
\end{solution}

\begin{solution}{exo:b2:reaction-free-energy:9}
الاستبدال: \ce{CH3CHBrCH3 + OH- -> CH3CH(OH)CH3 + Br-}؛ والحذف:
\ce{CH3CHBrCH3 + OH- -> CH2=CHCH3 + H2O + Br-}. والفرق بين طاقتي جيبس
القياسيتين للتفاعلين هو $\Delta(\Delta_r G^\circ) = 40 - T \times
0.090$، وهو سالب فوق $40\,000/90 \approx \qty{440}{K}$. فالحذف يصنع
جزيئًا زائدًا على ما يصنعه الاستبدال: إنتروبياه أكبر، وحد الإنتروبيا
يكبر بالتناسب مع $T$.
\end{solution}

\begin{solution}{exo:b2:reaction-free-energy:10}
لكل جزيء من الجزيئات $N_A$ توجهان: $W = 2^{N_A}$، ومنه $S =
k_B\ln 2^{N_A} = R\ln 2 = \qty{5.76}{J/(K.mol)}$. فالبلورة ليست تامة
الانتظام: عند درجة حرارة منخفضة تتجمد الجزيئات في توجهاتها العشوائية
ولا تستطيع بلوغ الترتيب المنتظم الوحيد الذي يفترضه المبدأ
الثالث.
\end{solution}

\begin{solution}{exo:b2:reaction-free-energy:11}
نكامل $\dd\Delta_r S^\circ/\dd T = c/T$: $\Delta_r S^\circ(T) = \Delta_r
S^\circ(T_0) + c\ln(T/T_0)$؛ ويعطي \omterm{thm:b2:reaction-enthalpy:kirchhoff}{قانون كيرشوف} $\Delta_r H^\circ(T) = \Delta_r
H^\circ(T_0) + c(T - T_0)$؛ ثم نطرح $T\Delta_r S^\circ(T)$. وللحجر الجيري عند
\qty{1100}{K}: إلينغهام $178.32 - 1100 \times 0.16064 = \qty{1.62}{kJ/mol}$؛
والتصحيح $c[T - T_0 - T\ln(T/T_0)] = -0.00195 \times (801.85 - 1436.3) =
\qty{+1.24}{kJ/mol}$: $\Delta_r G^\circ = \qty{2.86}{kJ/mol}$. والخطأ،
\qty{1.2}{kJ/mol}، أقل من واحد في المئة من $\Delta_r H^\circ$، لكنه قرب
\omterm{def:b2:reaction-free-energy:inversion-temperature}{درجة حرارة الانقلاب} يحسم الإشارة.
\end{solution}

\begin{solution}{exo:b2:reaction-free-energy:12}
(أ) $\Delta_r H^\circ = -763.2 + 944.0 = \qty{+180.8}{kJ/mol}$، $\Delta_r
S^\circ = 353.2 + 205.15 - 50.62 - 2(223.08) = \qty{+61.6}{J/(K.mol)}$؛
$\Delta_r G^\circ(1000) = 180.8 - 61.6 = \qty{+119.2}{kJ/mol}$: التفاعل \omterm{def:b2:reaction-free-energy:exergonic}{ماص
للطاقة}، فلا كلورة. (ب) من أجل \ce{C + O2 -> CO2}، $\Delta_r G^\circ(1000) =
-393.5 - 2.9 = \qty{-396.4}{kJ/mol}$. وللمجموع، \ce{TiO2(s) + 2Cl2(g) + C(s)
-> TiCl4(g) + CO2(g)}، $\Delta_r G^\circ = \qty{-277.2}{kJ/mol}$: فالكربون
يستهلك ثنائي الأكسجين الذي يصنعه التفاعل الأول، والتفاعل المقترن \omterm{def:b2:reaction-free-energy:exergonic}{طارد
للطاقة} بقوة.
\end{solution}

\begin{solution}{pb:b2:reaction-free-energy:1}
\textbf{1.} \ce{CaCO3(s) -> CaO(s) + CO2(g)}.
\textbf{2.} $\Delta_r H^\circ = -635.09 - 393.51 + 1206.92 =
\qty{+178.3}{kJ/mol}$: \omterm{def:b2:reaction-enthalpy:exothermic}{ماص للحرارة}.
\textbf{3.} $\Delta_r S^\circ = 39.75 + 213.79 - 92.9 = \qty{+160.6}{J/(K.mol)}$،
وهي موجبة لأن غازًا يتكوّن من جسم صلب.
\textbf{4.} $\Delta_r G^\circ = 178.32 - 298.15 \times 0.16064 =
\qty{+130.4}{kJ/mol}$.
\textbf{5.} نعم: يلزم $\Delta_r G < 0$ للتفكك، وهي هنا موجبة بقوة؛ وثنائي
أكسيد الكربون في الهواء، بضغط أدنى بكثير من \qty{1}{bar}،
يخفض $\Delta_r G$ لكن بأقل بكثير من \qty{130}{kJ/mol} (الفصل التالي).
\textbf{6.} $\Delta_r C^\circ_p = 42.80 + 37.13 - 81.88 =
\qty{-1.95}{J/(K.mol)}$، وهي ضئيلة: فلا يكاد $\Delta_r H^\circ$ و$\Delta_r S^\circ$
يتغيران مع $T$.
\textbf{7.} $\Delta_r G^\circ(T) = 178.32 - 0.16064\,T$ (\unit{kJ/mol}).
\textbf{8.} \qty{+17.7}{kJ/mol} عند \qty{1000}{K}؛ و\qty{-30.5}{kJ/mol} عند
\qty{1300}{K}.
\textbf{9.} $T_i = 178\,320/160.64 = \qty{1110}{K}$.
\textbf{10.} $\Delta_r G^\circ/T = 178.32/T - 0.16064$، ومشتقتها
$-178.32/T^2 = -\Delta_r H^\circ/T^2$.
\textbf{11.} عند $T_i$، $c[T_i - T_0 - T_i\ln(T_i/T_0)] = -0.00195 \times
(811.9 - 1459.4) = \qty{+1.26}{kJ/mol}$؛ فتنعدم $\Delta_r G^\circ$
عند درجة أعلى بمقدار $1260/160.6 \approx \qty{8}{K}$، أي نحو \qty{1118}{K}: أقل من
واحد في المئة.
\textbf{12.} عند \qty{1000}{K} يكون $\Delta_r G > 0$: لا يمكن أن يجري التفكك
(بل يمتص الجير ثنائي أكسيد الكربون). وعند \qty{1300}{K} يكون $\Delta_r G <
0$: فيتفكك الحجر الجيري.
\textbf{13.} تصير $\Delta_r G$ موجبة: فيمتص الجير ثنائي أكسيد الكربون
ويعود كربونات.
\textbf{14.} $\delta S_{\text{c}} = -\Delta_r G\,\dd\xi/T = 30\,506/1300 =
\qty{23.5}{J/K}$ لكل مول.
\textbf{15.} نعم: سيبيّن \cref{ch:b2:chemical-potential} أن
$\Delta_r G = \Delta_r G^\circ + RT\ln(p_{\ce{CO2}}/p^\circ)$، فتنخفض
حين يكون ضغط ثنائي أكسيد الكربون أدنى من \qty{1}{bar}، بحيث يبدأ
التفكك تحت $T_i$.
\textbf{16.} $n = 10^6/56.08 = \qty{1.783e4}{mol}$.
\textbf{17.} $1.783 \times 10^4 \times 178.32 = \qty{3.18e6}{kJ}$، أي
\qty{3.18}{GJ} لكل طن.
\textbf{18.} $1.783 \times 10^4 \times 44.01 = \qty{785}{kg}$ من \ce{CO2}.
\textbf{19.} الحرارة الواجب توفيرها $3.18 \times 10^6/0.50 = \qty{6.36e6}{kJ}$، ومنه
$6.36 \times 10^6/802.3 = \qty{7.93e3}{mol}$ من الميثان، أي \qty{127}{kg}.
\textbf{20.} $7.93 \times 10^3 \times 44.01 = \qty{349}{kg}$ من \ce{CO2}.
\textbf{21.} $785 + 349 = \qty{1134}{kg}$ لكل طن من الجير، 69\,\% منها
مصدره الحجر الجيري.
\textbf{22.} ثنائي أكسيد الكربون الآتي من الحجر الجيري تحدده الستوكيومترية،
أيًّا كان ما يسخّن الفرن؛ ولا يمكن خفض إلا حصة الوقود.
\textbf{23.} يتفكك الحجر الجيري تحت \qty{1}{bar} من ثنائي أكسيد الكربون فوق
درجة حرارة انقلابه، $\boldsymbol{T_i \approx \qty{1.11e3}{K}}$
(نحو \qty{840}{\celsius}).
\end{solution}
